% Conjectura manuscript template. Plain LaTeX with standard packages only, so
% that the source is accepted by arXiv without modification.
\documentclass[11pt]{amsart}
\usepackage{amsmath,amssymb,amsthm}
\usepackage[colorlinks=true,allcolors=blue]{hyperref}

\newtheorem{theorem}{Theorem}
\newtheorem{lemma}[theorem]{Lemma}
\newtheorem{proposition}[theorem]{Proposition}
\newtheorem{corollary}[theorem]{Corollary}
\theoremstyle{definition}
\newtheorem{definition}[theorem]{Definition}
\newtheorem{conjecture}[theorem]{Conjecture}
\theoremstyle{remark}
\newtheorem{remark}[theorem]{Remark}

\title{Title}
\author{Author}
\address{Universidade Federal de Goi\'as (UFG), Brazil}
\date{\today}

\begin{document}

\begin{abstract}
State the result exactly, with all hypotheses.
\end{abstract}

\maketitle

\section{Introduction}
The conjecture, its source and statement number, what was known, what is proved here.

\section{Statement and definitions}

\section{Proof}

\section{Computational evidence}
Exact range covered, software versions, and the commands that reproduce it.

\section{Formal verification}
Lean theorem names, toolchain and Mathlib revision, receipt identifier, and the
scope of what is and is not formalized.

\section*{Use of AI systems}
Which systems were used, for which parts of the work, and how each part was
verified. See the attribution record in the repository.

\section*{Limitations}

\begin{thebibliography}{9}
\end{thebibliography}

\end{document}
